% ICIP 2026 Grand Challenge — ClearSAR Track 1
% Two-page extended abstract (spconf.sty + IEEEbib.bst, per author kit)

\documentclass{article}
\usepackage{spconf,amsmath,amssymb,graphicx,booktabs}
\usepackage[hidelinks]{hyperref}
\usepackage{microtype}

\title{RFI DETECTION IN SENTINEL-1 SAR IMAGERY VIA PSEUDO-LABEL
DISTILLATION AND MULTI-RESOLUTION ENSEMBLE}

\name{Anas Elghoudane}
\address{Independent Researcher (\emph{UdeM AI} team)\\
\texttt{anaselghoudane@gmail.com}}

\begin{document}
\ninept
\maketitle

\begin{abstract}
We describe the \emph{UdeM AI} team's solution to the ESA ClearSAR
Track 1 Grand Challenge: Radio-Frequency Interference (RFI) detection
in Sentinel-1 SAR quicklooks. Our pipeline combines a YOLO26x base
detector with pixel-space Normalised Wasserstein Distance (NWD)
regression, pseudo-label distillation from a multi-source ensemble
teacher onto unlabelled test imagery, multi-resolution test-time
augmentation (TTA) fused by Weighted Box Fusion (WBF), and a selective
cross-fold ensemble that drops the domain-overlapping fold. Ablation
on a 401-image leak-isolated meta-validation set shows
pseudo-distillation contributes $+0.041$ mAP$_{50\text{-}95}$ over the
per-fold baseline, multi-resolution TTA adds $+0.034$, and selective
fold dropping adds $+0.002$, for a total $+0.077$ over the per-fold
meta-validation baseline. The final submission reaches 0.4776
mAP$_{50\text{-}95}$ on the held-out test set, placing 28th of 120
teams. All training and fine-tuning used a single cloud NVIDIA A100
GPU.
\end{abstract}

\begin{keywords}
Ensemble, pseudo-label distillation, radio-frequency interference,
SAR imagery, small object detection.
\end{keywords}

\section{TASK AND DATA}
\label{sec:task}

The ESA ClearSAR Track~1 Grand Challenge \cite{clearsar2026} frames
RFI detection in Sentinel-1 quicklooks as a COCO-style object
detection task with mAP averaged over IoU 0.50--0.95. The labelled
set contains 3{,}154 images; 786 test images are held out by the
organisers and scored by the challenge platform from a zipped
COCO-format JSON of detections. RFI appears as thin horizontal stripes (median height
$\approx$10\,px, median width $\approx$140\,px) in variable-size RGB
quicklooks (median $\approx$515$\times$342\,px), with a median
per-box aspect ratio of $\approx$8:1. Three constraints follow:
IoU-based regression gives no gradient for non-overlapping thin boxes;
SAR side-looking geometry breaks horizontal symmetry, so flip TTA is
disabled (H-flip costs $-6.5\%$ \cite{oliver2004sar}); and inference
resolutions above 640\,px amplify speckle relative to signal. We
reserve a leak-isolated meta-validation set of 401 images sampled
from the fold-0 validation pool (seed 42) for honest local
evaluation.

\section{METHOD}
\label{sec:method}

\begin{figure*}[t]
\centering
\setlength{\fboxsep}{3pt}
\scriptsize
\begin{tabular}{@{}c@{$\,\to\,$}c@{$\,\to\,$}c@{$\,\to\,$}c@{$\,\to\,$}c@{$\,\to\,$}c@{$\,\to\,$}c@{}}
\fbox{\parbox[c][12mm][c]{17mm}{\centering YOLO26x base\\5 folds\\CLAHE\,+\,NWD}} &
\fbox{\parbox[c][12mm][c]{15mm}{\centering V12 teacher\\WBF\\ensemble}} &
\fbox{\parbox[c][12mm][c]{17mm}{\centering Pseudo-label\\786 test imgs\\conf\,$\ge$\,0.75}} &
\fbox{\parbox[c][12mm][c]{17mm}{\centering Student FT\\$f\!\in\!\{1,2,4\}$\\freeze 10}} &
\fbox{\parbox[c][12mm][c]{18mm}{\centering Multi-res TTA\\608/640/\\672/704\,+\,WBF}} &
\fbox{\parbox[c][12mm][c]{17mm}{\centering Cross-fold WBF\\drop $f_0$\\$\to S^{*}$}} &
\fbox{\parbox[c][12mm][c]{19mm}{\centering Blend $w{=}0.25$\\$S^{*}\!\oplus$\,V12\\$\to$ V17 (0.4776)}} \\
\end{tabular}
\caption{Pipeline. A YOLO26x base detector (CLAHE input, pixel-space
NWD loss) is trained per fold; the V12 pre-distillation ensemble
pseudo-labels the 786 test images; per-fold students are fine-tuned
on the fold-augmented data, run multi-resolution TTA fused by WBF,
fused across folds (dropping the domain-overlapping fold~0) into a
super-source $S^{*}$, and blended with V12 at $w{=}0.25$.}
\label{fig:pipeline}
\end{figure*}

The pipeline (Figure~\ref{fig:pipeline}) has four stages.

\textbf{Base detector.} YOLO26x \cite{ultralytics2024yolo11}
($\approx$55.7\,M parameters, COCO pre-trained, no P2 head) trained
with a CIoU\,+\,NWD box loss
$\mathcal{L}_{\mathrm{box}} = (1{-}\alpha)\mathcal{L}_{\mathrm{CIoU}}
{+}\alpha\mathcal{L}_{\mathrm{NWD}}$ at $\alpha{=}0.5$, with NWD
\cite{wang2021nwd} computed in pixel space (normalisation constant
$C{=}12.8$\,px). Ultralytics' \texttt{BboxLoss.forward} receives
boxes in stride-relative grid units, so we multiply by the anchor
stride before computing NWD; this keeps gradients non-saturating for
the small heights characteristic of RFI ($\approx$10\,px median).
Images are pre-processed with CLAHE (clip 3.0, $8\times8$ tiles) on
the LAB L-channel. Horizontal and vertical flips are disabled.
Training uses MuSGD for up to 120 epochs (patience 20) at 640\,px,
batch 16, $\eta_0{=}10^{-3}$ with cosine decay to $10^{-5}$.

\textbf{Pseudo-label distillation.} The teacher (V12) is our
strongest pre-distillation ensemble: a single-stage WBF
\cite{solovyev2021wbf} over multiple YOLO26x fold and
multi-resolution-TTA sources together with one RF-DETR-Large
\cite{roboflow2024rfdetr} fold checkpoint at weight $1.5\times$,
fused with $\tau{=}0.70$ and \texttt{conf\_type}=max. V12 scores
0.4755 on the held-out test set. The teacher is applied to the 786
unlabelled test images with confidence threshold 0.75 (max 20
boxes/img), producing 600 pseudo-labelled images. For each fold
$f\in\{1,2,4\}$, a student is fine-tuned from the per-fold base on
the fold-augmented dataset with AdamW \cite{loshchilov2019adamw}
($\eta_0{=}10^{-4}$, \texttt{freeze}=10, 15 epochs, batch 24, same
NWD blend). Fold 0 is omitted because its base partition overlaps the
meta-validation pool.

\textbf{Multi-resolution TTA and fusion.} Each student is inferred at
$\mathcal{I}=\{608,640,672,704\}$\,px; WBF across resolutions per
fold yields one TTA source. The three per-fold TTA sources are then
fused with equal weights into a super-source $S^{*}$. The final
submission V17 is the two-source WBF blend $D_{\mathrm{final}} =
\mathrm{WBF}(D_{\mathrm{V12}}, S^{*}; w{=}(1.0, 0.25), \tau{=}0.70)$.

\section{RESULTS AND DISCUSSION}
\label{sec:results}

\begin{table}[t]
\centering
\caption{Component ablation on the 401-image meta-validation set.}
\label{tab:ablation}
\small
\begin{tabular}{lcc}
\toprule
\textbf{Configuration} & \textbf{mAP$_{50\text{-}95}$} & $\Delta$ \\
\midrule
Per-fold baseline (fold 0)            & 0.394 & --- \\
+ Pseudo-FT, fold 4 only              & 0.435 & $+0.041$ \\
+ Multi-res TTA, fold 4 only          & 0.469 & $+0.034$ \\
+ 3-fold (drop $f_0$), equal weights  & 0.471 & $+0.002$ \\
+ V12 blend @ $w{=}0.25$ (V17)        & 0.471$^\dagger$ & --- \\
\bottomrule
\end{tabular}\\[2pt]
{\scriptsize $^\dagger$Meta-val cannot exceed V12 by adding an honest
source (V12 is leaked on this split). Official held-out test:
\textbf{0.4776}.}
\end{table}

Table~\ref{tab:ablation} reports the component ablation.
Pseudo-distillation contributed the largest single gain ($+0.041$):
the soft signal in the teacher's confidence scores compensates for
the small training set more effectively than the architectural
alternatives we evaluated. Multi-resolution TTA added a further
$+0.034$ because thin stripes a few pixels tall are resolved unevenly
by YOLO's strided feature pyramid at any single input size. Dropping
the domain-overlapping fold-0 from the cross-fold ensemble was
net-positive: correlated predictions drag the WBF consensus when one
source has seen the validation data during training. Blending the
V12 teacher into the final fusion does not improve the leaked
meta-val baseline but transfers cleanly to held-out test
(V12 alone $= 0.4755 \to $ V17 $= 0.4776$).

\textbf{Negative results.} A second architectural family (RF-DETR)
plateaued at $+0.0002$ across six checkpoint--weight variants;
D-FINE-X \cite{peng2024dfine} fine-tuned on fold~0 for 10 epochs
reached $3\times10^{-4}$ mAP, suggesting DETR-family decoders with
different matching \cite{huang2024deim} would need substantially
more epochs to adapt to the elongated geometry of RFI; inference-only
SAHI \cite{akyon2022sahi} at 320-px tiles regressed $-0.085$,
fragmenting long horizontal stripes at tile boundaries; and
temperature-scaled calibration was at most $+0.0002$.

\section{CONCLUSION}
\label{sec:conclusion}

The final pipeline reaches 0.4776 mAP$_{50\text{-}95}$ on the
held-out test set; the meta-validation ablation chain
(Table~\ref{tab:ablation}) totals $+0.077$ over the per-fold
baseline. The submission placed the \emph{UdeM AI} team 28th of 120
teams on the public leaderboard of ESA ClearSAR Track~1. All training and fine-tuning were performed on
a single cloud NVIDIA A100 GPU. The largest remaining gap is at high
IoU thresholds, where sub-pixel height accuracy dominates and a
640-px input under-resolves the median $\approx$10-px stripe; sliced
\emph{training} at 320-px tiles is the most promising next move.

\bibliographystyle{IEEEbib}
\begin{thebibliography}{9}

\bibitem{clearsar2026}
European Space Agency,
``ClearSAR Track 1: Radio-frequency interference detection in
Sentinel-1 quicklook imagery,''
Challenge specification, 2026.

\bibitem{oliver2004sar}
C.~Oliver and S.~Quegan,
\emph{Understanding Synthetic Aperture Radar Images}.
SciTech Publishing, 2004.

\bibitem{ultralytics2024yolo11}
G.~Jocher and J.~Qiu,
``Ultralytics YOLO26,''
\url{https://github.com/ultralytics/ultralytics}, 2026.

\bibitem{wang2021nwd}
J.~Wang, C.~Xu, W.~Yang, and L.~Yu,
``A normalized Gaussian Wasserstein distance for tiny object
detection,''
\emph{arXiv:2110.13389}, 2021.

\bibitem{solovyev2021wbf}
R.~Solovyev, W.~Wang, and T.~Gabruseva,
``Weighted boxes fusion,''
\emph{Image and Vision Computing}, vol.~107, p.~104117, 2021.

\bibitem{roboflow2024rfdetr}
Roboflow,
``RF-DETR: A real-time transformer-based object detector,''
\url{https://github.com/roboflow/rf-detr}, 2024.

\bibitem{loshchilov2019adamw}
I.~Loshchilov and F.~Hutter,
``Decoupled weight decay regularization,''
in \emph{Proc. ICLR}, 2019.

\bibitem{peng2024dfine}
Y.~Peng, H.~Li, P.~Wu, Y.~Zhang, X.~Sun, and F.~Wu,
``D-FINE: Redefine regression task of DETR as fine-grained
distribution refinement,''
\emph{arXiv:2410.13842}, 2024.

\bibitem{huang2024deim}
S.~Huang, Z.~Lu, X.~Cun, Y.~Yu, X.~Zhou, and X.~Shen,
``DEIM: DETR with improved matching for fast convergence,''
\emph{arXiv:2412.04234}, 2024.

\bibitem{akyon2022sahi}
F.~C. Akyon, S.~O. Altinuc, and A.~Temizel,
``Slicing aided hyper inference and fine-tuning for small object
detection,''
in \emph{Proc. IEEE ICIP}, 2022, pp.~966--970.

\end{thebibliography}

\end{document}
